\section{Introduction}
\label{sec:introduction}

\subsection{Two Families of Structure-Layer Algorithms}

The redistricting platform's three-layer compositor separates plan generation
into Structure (how each bisection node is solved), WeightsOverride (what edge
weights encode), and SeedCompositor (how many seeds are tried).
The B-series has developed a rich family of Structure-layer algorithms,
all sharing a common pattern: they operate on the graph topology of the
census-tract adjacency and optimise some variant of the edge-cut objective.
\metis{} minimises raw edge cut via multilevel KL refinement~\citep{karypis1998};
\geosec{} (T.1) normalises edge cut by the isoperimetric ratio
$\sqrt{\min(i,k-i)}$ to favour balanced tree structures~\citep{b8paper};
\areasec{} (T.2) adds an area-swing constraint.
\sa{} (U.3) replaces the single \metis{} call with a simulated annealing
refinement loop to escape local optima in the edge-cut landscape~\citep{b19paper}.

All of these algorithms are, fundamentally, graph-topology methods: they
reason about which edges to cut, which vertices to separate, and how to
minimise the cost of a topological partition.

\paragraph{A second family: proximity packing.}
This paper introduces a different approach.
Centroidal Voronoi Districts (\cvd{}) assigns each census tract to its nearest
seed and iterates seed placement to convergence.
No edge-cut objective is consulted.
In graph-distance mode, the distance is BFS hop count on the tract adjacency
graph; in geographic mode, the distance is Euclidean distance on projected
tract centroids.
The common idea is \emph{proximity}: which seed owns each tract, and where
should that seed move after the ownership regions are known?
\cvd{} is therefore in the tradition of Voronoi-based spatial partitioning
rather than graph partitioning, while still using graph connectivity for
hop-count distance and contiguity repair.

The conceptual distinction matters in practice.
Edge-cut minimisation produces plans that follow existing graph structure:
METIS tends to find the ``thinest'' cuts in the adjacency graph, which often
align with roads, rivers, and county boundaries because those geographic
features create naturally low-weight edges (short shared boundaries).
\cvd{} instead fills geographic space radially from seed points.
In convex regions, \cvd{} produces roughly circular districts.
In elongated or irregular regions, it produces roughly elliptical districts
aligned with the long axis.
The resulting compactness profile is fundamentally different from the
compactness achieved by edge-cut minimisation.

\subsection{Lloyd's Algorithm and the CVD Connection}

\cvd{} is the discrete redistricting analogue of Lloyd's algorithm for
optimal quantisation~\citep{lloyd1982}.
Lloyd's algorithm, in its Euclidean form, iterates:
\begin{enumerate}
\item Assign each point to its nearest centre.
\item Move each centre to the centroid (population-weighted mean) of its cluster.
\item Repeat until convergence.
\end{enumerate}
Du, Faber and Gunzburger~\citeyear{du1999} proved convergence for Centroidal
Voronoi Tessellations (CVTs) under mild regularity conditions and surveyed their
applications in mesh generation, image compression, and spatial discretisation.

\cvd{} has two implemented metric modes:
\begin{itemize}
\item \textbf{Graph-distance CVD}: replaces Euclidean distance with BFS
  hop-count distance on the census-tract adjacency graph, and replaces the
  continuous centroid with a graph-distance medoid.
\item \textbf{Geographic CVD}: uses Euclidean distance on Albers-projected
  tract centroids, computes a population-weighted coordinate mean, and snaps
  the result back to the nearest actual tract.
\end{itemize}
Both modes keep seeds as real census tracts because the next Voronoi step needs
a valid tract or coordinate anchor that can be replayed from public inputs.
The graph-distance mode remains useful when centroid companion data are absent;
the geographic mode is the more literal Lloyd-style construction when centroid
data are present.

\subsection{Use Cases for CVD}

Three practical use cases motivate \cvd{}:

\paragraph{Fast initial plan.}
\cvd{} runs with no spanning trees, no MCMC, and no parameter tuning beyond
the iteration count and metric choice.
For exploratory runs or as a starting point for \sa{} refinement, \cvd{}
provides a valid, balanced, compact plan at low cost.

\paragraph{Geometric baseline.}
The ``most geometrically natural partition'' is best treated as a comparative
baseline in redistricting litigation, not as a binding legal rule.
\cvd{} provides this baseline: the plan produced by assigning each tract to
its nearest district centre, with no reference to political boundaries,
demographic composition, or incumbent addresses.
This is distinct from the \metis{} baseline, which minimises graph-topological
boundary length rather than geographic proximity.

\paragraph{Compactness diversity.}
The ensemble literature~\citep{deford2021} shows that plan quality is best
assessed by comparing a map against the space of alternatives.
Adding \cvd{} plans to the ensemble diversifies the compactness profile:
geometric proximity and edge-cut minimisation are complementary objectives
that explore different regions of the plan space.

\subsection{Main Results}

\paragraph{Algorithm.}
We define \cvd{} bisection (Section~\ref{sec:algorithm}) as a drop-in
Structure-layer replacement for \metis{}.
The algorithm has one primary tunable parameter:
\texttt{cvd\_iters} (default 20) controls the maximum number of
assignment-medoid iterations.
Seed initialisation uses farthest-first selection: the first seed is derived
deterministically from the global base seed and node path; the second is the
tract farthest from the first under the selected metric.
A post-hoc 200-iteration boundary-swap rebalance enforces the population
balance tolerance.

\paragraph{Implementation boundary.}
The current implementation is a binary splitter inside
\texttt{BisectionTree::from\_k}.  For an odd-seat node such as $k=7$, the
tree schedule records a $3:4$ child split, but the CVD rebalance currently
targets a half-population split at the local node.
This is acceptable for explaining the CVD mechanism and for even-split
fixtures, but publication claims for odd-seat statewide runs need either a
ratio-aware CVD rebalance or archived evidence showing the final RPLAN package
passes population audit after the full tree completes.

\paragraph{Convergence.}
For census-tract graphs, draft \cvd{} runs typically converge in 8--15
iterations on large states (NC, TX).
For path-like or elongated subgraphs (rural county chains), farthest-first seeding
trivially places seeds at opposite ends of the path, and convergence occurs in
1--3 steps.
The algorithm halts when seeds are unchanged between consecutive iterations.

\paragraph{Compactness profile.}
\cvd{} achieves Polsby-Popper compactness competitive with \metis{} in the
draft NC ($k=14$) comparison despite not optimising edge cuts.
The compactness arises naturally from the geometric proximity objective: tracts
near the centre of a district tend to be compact because they are equally
distant from all boundaries.
Section~\ref{sec:comparison} gives the three-algorithm comparison
(\cvd{}, \metis{}, \sa{}) on NC, WI, and TX, but treats the numbers as
benchmark targets until they are paired with archived run packages.

\paragraph{Geographic mode.}
Section~\ref{sec:phase2} describes Geographic \cvd{}: replacing BFS hop-count
distance with Euclidean distance on projected tract centroids.
The code now exposes \texttt{VoronoiMetric::Geographic} and loads optional
\texttt{tract\_centroids} companion data; production claims require real
centroid-backed packages rather than projection-only expectations.

\subsection{Paper Structure}

Section~\ref{sec:background} reviews Lloyd's algorithm, CVT theory, and
graph-distance medoid clustering.
Section~\ref{sec:algorithm} gives the formal \cvd{} algorithm with pseudocode,
the seed derivation, and convergence analysis.
Section~\ref{sec:comparison} presents empirical comparisons on NC, WI, and TX.
Section~\ref{sec:phase2} describes geographic \cvd{}.
Section~\ref{sec:conclusion} situates \cvd{} in the B-series algorithm landscape.

\subsection{Relation to Prior Work}

\cvd{} is a Structure-layer algorithm.
Seed-compositor strategies such as convergence sweeps and \be{} can run over
it without changing the split rule.
It is related to but distinct from \metis{}~\citep{karypis1998} (graph
topology vs.\ geometry) and from \sa{}~\citep{b19paper} (no warm start,
different objective).
It is the companion to \bfsgrowth{} (T.12): BFS growth from k seeds also
uses geographic proximity but grows greedily rather than iterating to
convergence.
Lloyd's algorithm in redistricting is not new as a concept, but prior
applications~\citep{du1999} use Euclidean distance and continuous coordinates;
the graph-distance medoid formulation for census-tract graphs is, to the
authors' knowledge, novel.
